% Thesis Abstract -----------------------------------------------------
%
% -------------------------------------------------------------------- 
%\pagenumbering{roman} \setcounter{page}{1} 
%\noindent{\large  \textbf{{\fontfamily{ptm}\selectfont Supergraph based Periodic Behavior Mining in Dynamic Networks}}}\\
%by \\
%Sajal Halder  


% \chapter*{Abstract\markboth{Abstract}{Abstract}}

% Hepatocellular carcinoma (HCC) is a highly prevalent and lethal malignancy with substantial molecular heterogeneity, making accurate disease-stage prediction and therapeutic target identification challenging. This study developed an integrated computational framework combining multi-omics analysis, machine learning, explainable artificial intelligence (XAI), molecular docking, and computational drug discovery to investigate HCC progression. RNA-sequencing, DNA methylation, and clinical data were obtained from the TCGA-LIHC dataset and analyzed to identify stage-associated molecular features. Machine learning models including Support Vector Machine (SVM), Random Forest (RF), XGBoost, and ensemble learning were developed for early- and advanced-stage classification. The ensemble model achieved the best performance, with a balanced accuracy of 74.5\% and ROC-AUC of 84\%. External validation using the independent GEO dataset GSE14520 confirmed eight biomarkers: ATOX1, PCCB, BICD2, EIF4E, CCL16, ANGPT2, VEGFA, and ECT2. Functional enrichment analysis indicated their involvement in angiogenesis, tumor proliferation, cellular growth regulation, and cancer-related signaling pathways. SHAP analysis further identified influential biomarkers contributing to model predictions and improved the interpretability of the developed model. For therapeutic investigation, PCCB, VEGFA and SPRY4 were selected as representative targets from RNA expression and DNA methylation analyses, respectively. Molecular docking using CB-Dock2 was performed to evaluate selected FDA-approved anticancer drugs against these targets. Drug-likeness, pharmacokinetic properties, and toxicity were subsequently assessed using SwissADME and ProTox-3.0. The docking and pharmacological analyses identified several candidate compounds with favorable predicted binding affinities, drug-like properties, and comparatively acceptable toxicity profiles. Overall, the proposed framework integrates selected biomarker discovery, predictive modeling, model interpretation, external validation, and therapeutic candidate screening into a unified computational workflow for prioritizing anticancer drugs investigation and stage progression in HCC. This approach may support future research on stage-specific HCC characterization and the development of potential biomarker-driven therapeutic strategies.
    
% \par \textit{\textbf{Keywords:}} Multiomics, TCGA, GEO, HCC, Biomarker, SHAP Analysis, Molecular Docking, Drug Discovery.



\chapter*{Abstract\markboth{Abstract}{Abstract}}
\addcontentsline{toc}{chapter}{Abstract}

% Bangladesh is highly vulnerable to flooding and water-logging due to its extensive river network, seasonal rainfall, and substantial hydrological variability across different regions. Accurate and timely water-level forecasting is therefore important for improving flood preparedness and early warning. However, conventional centralized forecasting approaches may not adequately capture the heterogeneous characteristics of different monitoring stations and may require the centralization of distributed data. This study proposes a Clustered Federated Learning--based Encoder--Decoder Long Short-Term Memory (CFL-EDLSTM) framework for multi-horizon water-level forecasting in Bangladesh. The study utilizes hydrological and meteorological observations from 26 selected monitoring stations within the Brahmaputra River system in Bangladesh and considers water level as the target variable, with discharge, rainfall, temperature, and evaporation used as input variables. To account for differences in hydrological and geographical characteristics, the monitoring stations are organized into four clusters: Upstream, Middle, Downstream, and Urban. Federated model training is then performed within these clusters, enabling the proposed framework to learn region-specific temporal patterns without requiring direct centralization of local training data. The forecasting performance of the proposed CFL-EDLSTM model will be evaluated against Vanilla LSTM , GRU, XGBoost, and Temporal Convolutional Network (TCN) baseline models. The framework will generate water-level forecasts for multiple lead times of 3 hours, 6 hours, 12 hours, 1 day, 3 days, 7 days, and 15 days, allowing performance to be assessed across both short- and long-term forecasting horizons. In addition, SHAP-based analysis will be used to improve model interpretability by examining the contribution of the input variables to the predicted water levels. To demonstrate the practical usability of the proposed framework, a web-based visualization interface will be developed with an interactive map of Bangladesh showing the 25 monitoring stations. Users will be able to select individual station markers and view current, previous, and predicted water-level values for different forecasting horizons. A dedicated comparison tab will also compare the water levels predicted by the proposed model with corresponding water-level information from the Flood Forecasting and Warning Centre (FFWC) of the Bangladesh Water Development Board, used as an external point of comparison rather than as validated ground truth. Overall, the proposed CFL-EDLSTM framework aims to provide an interpretable and privacy-preserving approach for multi-horizon water-level forecasting while accounting for regional heterogeneity through four distinct station clusters and demonstrating the results through an accessible visualization and comparison platform.

Bangladesh's extensive river network, seasonal rainfall, and considerable hydrological variability across regions leave the country highly exposed to flooding and water-logging, making accurate, timely water-level forecasting central to effective flood preparedness and early warning. Conventional centralized forecasting approaches, however, often struggle to capture the heterogeneous characteristics of different monitoring stations, and typically require pooling data from geographically dispersed sources into a single location. This thesis proposes a Clustered Federated Learning-based Encoder–Decoder Long Short-Term Memory (CFL-EDLSTM) framework for multi-horizon water-level forecasting that addresses both limitations at once.

The study draws on hydrological and meteorological observations from 25 monitoring stations along the Brahmaputra river system in Bangladesh, using water level as the target variable and discharge, rainfall, temperature, and evaporation as input variables. To account for differences in hydrological and geographical setting, the stations are grouped into four clusters — Upstream, Middle, Downstream, and Urban — and federated training is carried out within each cluster, allowing the framework to learn region-specific temporal patterns without centralizing any station's raw data.

CFL-EDLSTM's forecasting performance is evaluated against four baseline models — Vanilla LSTM, GRU, XGBoost, and a Temporal Convolutional Network (TCN) — across seven forecast horizons ranging from 3 hours to 15 days, allowing performance to be assessed across both short- and long-term forecasting behaviour. SHAP-based analysis is further used to examine how each input variable contributes to the model's predictions, adding a layer of interpretability that is often missing from deep learning-based hydrological forecasting.

To demonstrate the practical usability of the proposed framework, a web-based visualization interface presents an interactive map of the 25 monitoring stations, letting users select individual stations to view current, historical, and predicted water levels across the different forecast horizons. A dedicated comparison view sets these predictions against corresponding data from the Flood Forecasting and Warning Centre (FFWC) of the Bangladesh Water Development Board, used here as an external point of reference rather than as validated ground truth.

Taken together, this thesis aims to demonstrate that clustered federated learning can support accurate, interpretable, and privacy-preserving multi-horizon water-level forecasting while remaining sensitive to the regional heterogeneity of Bangladesh's river network, and to make that forecasting accessible through a working visualization and comparison platform.

% \vspace{1em}
\par \textit{\textbf{Keywords:}} Clustered Federated Learning; Encoder--Decoder LSTM; Water-Level Forecasting; Bangladesh; Multi-Horizon Forecasting; SHAP; Explainable AI